\documentclass[journal]{IEEEtran}

\usepackage[noadjust]{cite}




% *** GRAPHICS RELATED PACKAGES ***
%
\ifCLASSINFOpdf
   \usepackage[pdftex]{graphicx}
\usepackage{mathrsfs}
  % declare the path(s) where your graphic files are
   \graphicspath{{img/}}
  % and their extensions so you won't have to specify these with
  % every instance of \includegraphics
   \DeclareGraphicsExtensions{.pdf,.jpeg,.png}
\else
  % or other class option (dvipsone, dvipdf, if not using dvips). graphicx
  % will default to the driver specified in the system graphics.cfg if no
  % driver is specified.
  % \usepackage[dvips]{graphicx}
  % declare the path(s) where your graphic files are
  % \graphicspath{{../eps/}}
  % and their extensions so you won't have to specify these with
  % every instance of \includegraphics
  % \DeclareGraphicsExtensions{.eps}
\fi
% graphicx was written by David Carlisle and Sebastian Rahtz. It is
% required if you want graphics, photos, etc. graphicx.sty is already
% installed on most LaTeX systems. The latest version and documentation
% can be obtained at: 
% http://www.ctan.org/pkg/graphicx
% Another good source of documentation is "Using Imported Graphics in
% LaTeX2e" by Keith Reckdahl which can be found at:
% http://www.ctan.org/pkg/epslatex
%
% latex, and pdflatex in dvi mode, support graphics in encapsulated
% postscript (.eps) format. pdflatex in pdf mode supports graphics
% in .pdf, .jpeg, .png and .mps (metapost) formats. Users should ensure
% that all non-photo figures use a vector format (.eps, .pdf, .mps) and
% not a bitmapped formats (.jpeg, .png). The IEEE frowns on bitmapped formats
% which can result in "jaggedy"/blurry rendering of lines and letters as
% well as large increases in file sizes.
%
% You can find documentation about the pdfTeX application at:
% http://www.tug.org/applications/pdftex

\usepackage{algorithm}
\usepackage{algpseudocode}
\usepackage{booktabs} 
\makeatletter
\renewcommand\fs@ruled{%
  \def\@fs@cfont{\bfseries\color{blue}}%
  \let\@fs@capt\floatc@ruled
  \def\@fs@pre{{\color{blue}\hrule height.8pt depth0pt}\kern2pt}%
  \def\@fs@post{\kern2pt{\color{blue}\hrule}\relax}%
  \def\@fs@mid{\kern2pt{\color{blue}\hrule}\kern2pt}%
  \let\@fs@iftopcapt\iftrue
}
\makeatother

% *** Edit Packages ***
\usepackage{soul, color}
\usepackage{colortbl}
\usepackage{multirow}

\usepackage{enumitem}


% *** Placeholder ***
\usepackage{lipsum}


% *** MATH PACKAGES ***
%
\usepackage{amsmath}
\usepackage{amssymb}
\usepackage{amsfonts, bm}
\usepackage{steinmetz}
\usepackage{dsfont}
\usepackage[colorlinks=true]{hyperref} 
\hypersetup{
    colorlinks,
    citecolor=blue,
    linkcolor=blue,
    urlcolor=blue
}
\usepackage{xurl}      
\usepackage[caption=false, font=footnotesize]{subfig}
\usepackage{nomencl}
\makenomenclature

\usepackage{etoolbox}

\renewcommand\nomgroup[1]{%
  \item[\bfseries
  \ifstrequal{#1}{A}{General Math and Indices}{%
  \ifstrequal{#1}{B}{Grid Control Related}{%
  \ifstrequal{#1}{C}{RL Related}{}}}%
]}


% correct bad hyphenation here
\hyphenation{op-tical net-works semi-conduc-tor}


\makeatletter
\def\ps@IEEEtitlepagestyle{
  \def\@oddfoot{\mycopyrightnotice}
  \def\@evenfoot{}
}
\def\mycopyrightnotice{
  {\scriptsize
  \begin{minipage}{\textwidth}
  \end{minipage}
  }
}
\usepackage{booktabs,tabularx,makecell,array}


\usepackage{titlesec}

\titlespacing\section{0pt}{6pt plus 4pt minus 2pt}{0pt plus 2pt minus 2pt}
\titlespacing\subsection{0pt}{3pt plus 2pt minus 2pt}{0pt plus 2pt minus 2pt}
\titlespacing\subsubsection{0pt}{2pt plus 2pt minus 2pt}{0pt plus 2pt minus 2pt}

\usepackage{titlesec}
\usepackage{amsmath,amssymb,amsthm,epsfig,color,empheq,graphicx}

\newtheorem{proposition}{Proposition}
\newtheorem{lemma}{Lemma}
\newtheorem{corollary}{Corollary}
\newtheorem{theorem}{Theorem}
\newtheorem{definition}{Definition}
\newtheorem{remark}{Remark}
\newtheorem{example}{Example}

\DeclareMathOperator*{\argmax}{\arg\!\max}

\newcounter{assump}
\renewcommand{\theassump}{A.\arabic{assump}}

\newenvironment{assumption}{%
  \refstepcounter{assump}%
  \noindent\textbf{\textit{[\theassump]}}\;
}{\par}

\begin{document}
 
\title{Carbon-Cap-Constrained Planning of Shared Energy Storage\\ in a Microgrid With Batch-Resolved Carbon Accounting}
 
\author{Author~Name%
\thanks{Manuscript submitted \today. (\textit{Corresponding author: X. Y.})}%
\thanks{The authors are with the School of Electrical Engineering,
Southeast University, Nanjing, China (e-mail: author@seu.edu.cn).}}
 
\markboth{IEEE Transactions on Smart Grid,~Vol.~XX, No.~X, Month~2026}%
{Author \MakeLowercase{\textit{et al.}}: Planning of Shared Energy Storage in a Microgrid}
 
\maketitle
 
% =====================================================================
\section*{Nomenclature}
 
\subsubsection*{Sets and Indices}
\begin{IEEEdescription}[\IEEEsetlabelwidth{$\mathcal{N}^{\mathrm{es}}$}]
\item[$\mathcal{N}$] Set of buses, indexed by $i,j$; bus $0$ is the substation.
\item[$\mathcal{L}$] Set of distribution lines, indexed by $(i,j)$.
\item[$\mathcal{N}^{\mathrm{es}}$] Set of candidate shared-ESS buses.
\item[$\mathcal{G}_i$] Set of dispatchable generators at bus $i$.
\item[$\mathcal{S}$] Set of representative and extreme scenarios, indexed by $s$.
\item[$\mathcal{D},\mathcal{T}$] Sets of days and hourly intervals, indexed by $d,t$.
\item[$\mathcal{K}$] Set of charging vintages (virtual batches), indexed by $k$.
\end{IEEEdescription}
 
\subsubsection*{Parameters}
\begin{IEEEdescription}[\IEEEsetlabelwidth{$\overline{\mathscr{w}}_i$}]
\item[$\rho_s$] Annual occurrence factor of scenario $s$.
\item[$\Delta t$] Duration of one operating interval.
\item[$r_{ij},x_{ij}$] Resistance and reactance of branch $(i,j)$.
\item[$\eta_i^{\mathrm{ch}},\eta_i^{\mathrm{dc}}$] ESS charging/discharging efficiencies.
\item[$\underline{\mathrm{SOC}}_i,\overline{\mathrm{SOC}}_i$] ESS state-of-charge limits.
\item[$\overline{S}_{ij}$] Apparent-power limit of branch $(i,j)$.
\item[$\overline{\mathscr{w}}_i$] Carbon-intensity cap at bus $i$.
\item[$\mathrm{CRF}$] Capital recovery factor.
\end{IEEEdescription}
 
\subsubsection*{Planning Variables}
\begin{IEEEdescription}[\IEEEsetlabelwidth{$E_i^{\mathrm{es,cap}}$}]
\item[$u_i^{\mathrm{es}}$] Binary shared-ESS siting variable.
\item[$P_i^{\mathrm{es,cap}}$] Installed shared-ESS power capacity.
\item[$E_i^{\mathrm{es,cap}}$] Installed shared-ESS energy capacity.
\end{IEEEdescription}
 
\subsubsection*{Operating Variables}
\begin{IEEEdescription}[\IEEEsetlabelwidth{$P_{i,s,d,t}^{\mathrm{es,ch}}$}]
\item[$P_{i,s,d,t}^{\mathrm{es,ch}}$] Aggregate ESS charging power.
\item[$P_{i,s,d,t}^{\mathrm{es,dc}}$] Aggregate ESS discharging power.
\item[$e_{i,s,d,t}^{\mathrm{es}}$] Aggregate stored energy.
\item[$P_{ij,s,d,t},Q_{ij,s,d,t}$] Active/reactive branch power flow.
\item[$v_{i,s,d,t}$] Squared bus-voltage magnitude.
\item[$P_{ij,s,d,t}^{+},P_{ij,s,d,t}^{-}$] Forward/reverse active line flow.
\item[$\mathscr{w}_{i,s,d,t}$] Nodal carbon intensity.
\end{IEEEdescription}
 
% =====================================================================
% Model section only.
% Required packages in the main IEEE manuscript:
% \usepackage{amsmath,amssymb,bm,mathrsfs}

\raggedbottom

\section{Introduction}

\section{Shared Storage Planning With Batch-Resolved Carbon Accounting}
\label{sec:model}
Consider a grid-connected microgrid represented by a fixed radial graph
$\mathscr{G}=(\mathcal{N},\mathcal{L})$, where bus $0$ is the point of
common coupling with the upstream grid. The microgrid contains
dispatchable generators, photovoltaic (PV) units, conventional loads, and
a data center at bus $i^{\mathrm{DC}}$. Shared energy storage can be
installed at candidate buses
$\mathcal{N}^{\mathrm{es}}\subseteq\mathcal{N}$ and is centrally
dispatched to support the entire microgrid.

The microgrid operator determines the siting and rated power/energy
capacities of the shared storage facilities and jointly schedules them
with the other resources. The planning decisions are common to
all scenarios, whereas dispatch and workload schedules are scenario
dependent. The objective is to minimize annualized cost subject to the
operating constraints and nodal carbon-intensity limits. To retain the
carbon attributes of stored electricity, we develop a
\emph{charging-time-resolved virtual carbon inventory} model. Each charging
interval creates a virtual energy batch, termed a \emph{carbon vintage},
whose charging-time carbon intensity is preserved until discharge. This
retains the temporal heterogeneity of stored energy and makes its
carbon-compliance value depend on the vintages available for discharge.

\subsection{Objective}
\emph{1) Primary Economic Objective:}
Define the shared-storage planning variables as
\begin{equation}
\boldsymbol{x}^{\mathrm{pl}}
=
\{u_i^{\mathrm{es}},P_i^{\mathrm{es,cap}},E_i^{\mathrm{es,cap}}\}_{i\in\mathcal{N}^{\mathrm{es}}},
\label{eq:planning_variables}
\end{equation}
and let $\boldsymbol{y}^{\mathrm{op}}$ collect all scenario-dependent
variables introduced below. The primary problem is
\begin{equation}
\min_{\boldsymbol{x}^{\mathrm{pl}},\,\boldsymbol{y}^{\mathrm{op}}}
\quad
C^{\mathrm{ann}},
\label{eq:stage1_objective}
\end{equation}
where
\begin{equation}
C^{\mathrm{ann}}
=
C^{\mathrm{inv}}+C^{\mathrm{op}}.
\label{eq:annual_cost}
\end{equation}

The operating scenarios are
\begin{equation}
\mathcal{S}=\mathcal{S}^{\mathrm{yr}}\cup\{\mathrm{ex}\},\qquad
\mathcal{S}^{\mathrm{yr}}=\{\mathrm{sp},\mathrm{su},\mathrm{au},\mathrm{wi}\},
\label{eq:scenario_set}
\end{equation}
where $\mathrm{ex}$ is an extreme operating block. Each scenario contains
$D^{\mathrm{blk}}=5$ days with $T=24$ hourly intervals. Here,
$d\in\mathcal{D}=\{1,\ldots,D^{\mathrm{blk}}\}$,
$t\in\mathcal{T}=\{1,\ldots,T\}$, and $q(d,t)=(d-1)T+t$ maps the block
into a 120-hour sequence. If $N_s$ is the number of calendar days
represented by scenario $s$, then
\begin{align}
\sum_{s\in\mathcal{S}} N_s
&=365,
\label{eq:annual_days}\\
\rho_s
&=\frac{N_s}{D^{\mathrm{blk}}}.
\label{eq:scenario_factor}
\end{align}
Thus, each block is counted $\rho_s$ times; a synthetic extreme scenario
uses $N_{\mathrm{ex}}=0$ and affects feasibility only.

The annualized shared-storage investment cost is
\begin{equation}
C^{\mathrm{inv}}
=
\mathrm{CRF}
\sum_{i\in\mathcal{N}^{\mathrm{es}}}
\left(
c_i^{\mathrm{site}}u_i^{\mathrm{es}}
+
c_i^P P_i^{\mathrm{es,cap}}
+
c_i^E E_i^{\mathrm{es,cap}}
\right).
\label{eq:investment_cost}
\end{equation}
The annual operating cost is
\begin{equation}
C^{\mathrm{op}}
=c^{\mathrm{dem}}P^{\mathrm{peak}}
+\sum_{s\in\mathcal{S}}\rho_s
\sum_{d\in\mathcal{D}}\sum_{t\in\mathcal{T}}
\Gamma_{s,d,t}\Delta t,
\label{eq:operating_cost}
\end{equation}
where
\begin{align}
\Gamma_{s,d,t}
=\;&c_{s,d,t}^{\mathrm{grid}}P_{s,d,t}^{\mathrm{grid}}
+\sum_{i\in\mathcal{N}}\sum_{g\in\mathcal{G}_i}
c_g^{\mathrm{gen}}P_{g,s,d,t}
\nonumber\\
&+\sum_{i\in\mathcal{N}^{\mathrm{es}}}c_i^{\mathrm{deg}}
\left(P_{i,s,d,t}^{\mathrm{es,ch}}+P_{i,s,d,t}^{\mathrm{es,dc}}\right)
+c^{\mathrm{delay}}B_{s,d,t}
\nonumber\\
&+\sum_{i\in\mathcal{N}}c^{\mathrm{curt}}P_{i,s,d,t}^{\mathrm{curt}}
+c^{\mathrm{VOLL}}P_{s,d,t}^{\mathrm{shed,DC}}.
\label{eq:interval_cost}
\end{align}
Equations \eqref{eq:annual_cost}--\eqref{eq:interval_cost} account for
shared-storage investment, grid purchases, local generation, demand charges,
degradation, workload delay, PV curtailment, and data-center load shedding.

\emph{2) Carbon-Vintage Preservation Objective:}

Let $C^{\mathrm{op},*}$,
$u_i^{\mathrm{es},*}$, $P_i^{\mathrm{es,cap},*}$, and
$E_i^{\mathrm{es,cap},*}$ denote the optimal first-stage values. The
second stage fixes the storage planning decisions and optimizes only the
scenario-dependent operating variables:
\begin{equation}
\min_{\boldsymbol{y}^{\mathrm{op}}}
\quad
J^{\mathrm{res}},
\label{eq:stage2_objective}
\end{equation}
where
\begin{align}
J^{\mathrm{res}}
=
&
\sum_{s\in\mathcal{S}}
\rho_s
\sum_{d\in\mathcal{D}}
\sum_{t\in\mathcal{T}}
\sum_{i\in\mathcal{N}^{\mathrm{es}}}
\sum_{k\in\mathcal{K}}
\nonumber\\
&
\left(
\mathscr{w}^{\max}
-
\mathscr{w}_{i,k,s}^{\mathrm{es}}
\right)
P_{i,k,s,d,t}^{\mathrm{es,dc}}
\Delta t.
\label{eq:reserve_objective}
\end{align}
The corresponding restrictions are
\begin{align}
u_i^{\mathrm{es}}
&=
u_i^{\mathrm{es},*},
\label{eq:fix_siting}\\
P_i^{\mathrm{es,cap}}
&=
P_i^{\mathrm{es,cap},*},
\label{eq:fix_power_capacity}\\
E_i^{\mathrm{es,cap}}
&=
E_i^{\mathrm{es,cap},*},
\label{eq:fix_energy_capacity}\\
C^{\mathrm{op}}
&\leq
C^{\mathrm{op},*}
+
\varepsilon_C^{\mathrm{num}}.
\label{eq:economic_equivalence}
\end{align}

Because lower-carbon vintages receive larger coefficients in
\eqref{eq:reserve_objective}, minimizing $J^{\mathrm{res}}$ avoids their
unnecessary discharge. The tolerance $\varepsilon_C^{\mathrm{num}}$ only
absorbs solver precision and cannot alter the planning decisions. The
vintage-specific discharge variables define a virtual carbon-allocation
policy; they do not imply physical separation of energy inside the ESS.

\subsection{Constraints}
Unless otherwise stated, the following constraints apply to all
$s\in\mathcal{S}$, $d\in\mathcal{D}$, $t\in\mathcal{T}$, and all
applicable buses, lines, generators, and storage layers.

\emph{1) Data-Center Flexibility:}
The workload of the data center at bus $i^{\mathrm{DC}}$ is scheduled
over the full five-day chronology.

The backlog dynamics and inter-day transition are
\begin{align}
B_{s,d,t+1}
&=
B_{s,d,t}
+
A_{s,d,t}
-
x_{s,d,t},
\qquad
t\in\mathcal{T},
\label{eq:backlog_dynamic}\\
B_{s,d+1,1}
&=
B_{s,d,T+1},
\qquad
d=1,\ldots,D^{\mathrm{blk}}-1.
\label{eq:backlog_interday}
\end{align}
The boundary and capacity constraints are
\begin{equation}
B_{s,1,1}
=
0,
\qquad
B_{s,D^{\mathrm{blk}},T+1}
=
0,
\label{eq:backlog_boundary}
\end{equation}
\begin{equation}
0
\leq
x_{s,d,t}
\leq
\overline{x},
\qquad
0
\leq
B_{s,d,t}
\leq
\overline{B}.
\label{eq:workload_limits}
\end{equation}

Define the set of intervals no later than $(d,t)$ as
\begin{equation}
\mathcal{P}_{d,t}
=
\{
(d',t'):
q(d',t')
\leq
q(d,t)
\}.
\label{eq:past_intervals}
\end{equation}
The cumulative processed workload cannot exceed the cumulative arrivals:
\begin{equation}
\sum_{(d',t')\in\mathcal{P}_{d,t}}
x_{s,d',t'}
\leq
\sum_{(d',t')\in\mathcal{P}_{d,t}}
A_{s,d',t'}.
\label{eq:no_early_processing}
\end{equation}

For a maximum allowable delay of $L^{\max}$ intervals, define
\begin{equation}
\mathcal{H}_{d,t}
=
\{
(d',t'):
q(d',t')
\leq
\min[
q(d,t)+L^{\max},
D^{\mathrm{blk}}T
]
\}.
\label{eq:deadline_intervals}
\end{equation}
The deadline requirement is
\begin{equation}
\sum_{(d',t')\in\mathcal{H}_{d,t}}
x_{s,d',t'}
\geq
\sum_{(d',t')\in\mathcal{P}_{d,t}}
A_{s,d',t'}.
\label{eq:deadline_constraint}
\end{equation}

The IT power and total data-center power are
\begin{align}
P_{s,d,t}^{\mathrm{IT}}
&=
P^{\mathrm{idle}}
+
\kappa^{\mathrm{task}}
x_{s,d,t},
\label{eq:IT_power}\\
P_{s,d,t}^{\mathrm{DC}}
&=
P_{s,d,t}^{\mathrm{aux}}
+
\mathrm{PUE}_{s,d,t}
P_{s,d,t}^{\mathrm{IT}}.
\label{eq:DC_power}
\end{align}
Define the fixed-location indicator
\begin{equation}
\delta_i^{\mathrm{DC}}
=
\begin{cases}
1, & i=i^{\mathrm{DC}},\\
0, & i\neq i^{\mathrm{DC}}.
\end{cases}
\label{eq:DC_location_indicator}
\end{equation}
The active load at bus $i$ is then
\begin{equation}
P_{i,s,d,t}^{L}
=
P_{i,s,d,t}^{\mathrm{base}}
+
\delta_i^{\mathrm{DC}}
\left(
P_{s,d,t}^{\mathrm{DC}}
-
P_{s,d,t}^{\mathrm{shed,DC}}
\right).
\label{eq:nodal_load}
\end{equation}
The data-center load-shedding limit is
\begin{equation}
0
\leq
P_{s,d,t}^{\mathrm{shed,DC}}
\leq
P_{s,d,t}^{\mathrm{DC}}
-
P_{s,d,t}^{\mathrm{DC,crit}}.
\label{eq:DC_shedding_limit}
\end{equation}

Here, $A_{s,d,t}$, $x_{s,d,t}$, and $B_{s,d,t}$ denote arrived,
processed, and backlogged workload, respectively. Parameters
$\overline{x}$, $\overline{B}$, and $L^{\max}$ impose processing,
backlog, and delay limits.

\emph{2) Shared-Storage Planning and Aggregate Operation:}

The storage planning constraints are
\begin{align}
\underline{P}_i^{\mathrm{es,cap}}
u_i^{\mathrm{es}}
&\leq
P_i^{\mathrm{es,cap}}
\leq
\overline{P}_i^{\mathrm{es,cap}}
u_i^{\mathrm{es}},
\label{eq:storage_power_capacity}\\
\underline{E}_i^{\mathrm{es,cap}}
u_i^{\mathrm{es}}
&\leq
E_i^{\mathrm{es,cap}}
\leq
\overline{E}_i^{\mathrm{es,cap}}
u_i^{\mathrm{es}},
\label{eq:storage_energy_capacity}\\
\sum_{i\in\mathcal{N}^{\mathrm{es}}}
u_i^{\mathrm{es}}
&\leq
\overline{N}^{\mathrm{es}},
\label{eq:storage_site_limit}\\
h_i^{\min}P_i^{\mathrm{es,cap}}
&\leq
E_i^{\mathrm{es,cap}}
\leq
h_i^{\max}P_i^{\mathrm{es,cap}}.
\label{eq:storage_duration}
\end{align}

The aggregate charging power, discharging power, and stored energy are
obtained by summing all layers:
\begin{align}
P_{i,s,d,t}^{\mathrm{es,ch}}
&=
\sum_{k\in\mathcal{K}}
P_{i,k,s,d,t}^{\mathrm{es,ch}},
\label{eq:aggregate_charge}\\
P_{i,s,d,t}^{\mathrm{es,dc}}
&=
\sum_{k\in\mathcal{K}}
P_{i,k,s,d,t}^{\mathrm{es,dc}},
\label{eq:aggregate_discharge}\\
e_{i,s,d,t}^{\mathrm{es}}
&=
\sum_{k\in\mathcal{K}}
e_{i,k,s,d,t}^{\mathrm{es}}.
\label{eq:aggregate_energy}
\end{align}
The aggregate operating limits are
\begin{align}
0
\leq
P_{i,s,d,t}^{\mathrm{es,ch}}
&\leq
P_i^{\mathrm{es,cap}},
\label{eq:aggregate_charge_limit}\\
0
\leq
P_{i,s,d,t}^{\mathrm{es,dc}}
&\leq
P_i^{\mathrm{es,cap}},
\label{eq:aggregate_discharge_limit}\\
u_{i,s,d,t}^{\mathrm{es,ch}}
+
u_{i,s,d,t}^{\mathrm{es,dc}}
&\leq
1,
\label{eq:charge_discharge_exclusion}\\
P_{i,s,d,t}^{\mathrm{es,ch}}
&\leq
\overline{P}_i^{\mathrm{es,cap}}
u_{i,s,d,t}^{\mathrm{es,ch}},
\label{eq:charge_status}\\
P_{i,s,d,t}^{\mathrm{es,dc}}
&\leq
\overline{P}_i^{\mathrm{es,cap}}
u_{i,s,d,t}^{\mathrm{es,dc}},
\label{eq:discharge_status}\\
\underline{\mathrm{SOC}}_i
E_i^{\mathrm{es,cap}}
&\leq
e_{i,s,d,t}^{\mathrm{es}}
\leq
\overline{\mathrm{SOC}}_i
E_i^{\mathrm{es,cap}}.
\label{eq:SOC_limit}
\end{align}

The variables $u_i^{\mathrm{es}}$, $P_i^{\mathrm{es,cap}}$, and
$E_i^{\mathrm{es,cap}}$ are common to all scenarios. Consequently, the
seasonal and extreme scenarios jointly determine the final storage
location and capacity. The layer representation changes only the carbon
accounting resolution and does not alter the aggregate physical storage
constraints.

\emph{3) Charging-Time-Resolved Virtual Carbon Inventory:}
The water-tank carbon-footprint model represents stored energy and virtual
carbon emissions as a single well-mixed inventory \cite{chen2025copf}.
In contrast, the proposed model retains a separate virtual inventory for
each charging vintage. The associated energy and carbon remain
identifiable until discharge, allowing the planner to distinguish
low- and high-carbon stored energy and to optimize their discharge
allocation under nodal carbon-intensity limits.

\textit{a) Carbon-vintage creation:}

Within each representative five-day scenario, every hourly interval is
associated with one candidate charging-batch layer. Let $k$ denote the
layer index and define
\begin{equation}
\mathcal{K}
=
\{0,1,\ldots,K^{\max}\},
\qquad
K^{\max}
=
D^{\mathrm{blk}}T
=
120.
\label{eq:K_set}
\end{equation}
Layer $k=0$ represents the inventory that exists before the scenario
begins, whereas each layer $k\geq1$ corresponds to one candidate charging
interval. The formation day and hourly interval associated with layer
$k\geq1$ are
\begin{equation}
d_k
=
\left\lceil
\frac{k}{T}
\right\rceil,
\qquad
t_k
=
k-(d_k-1)T.
\label{eq:layer_mapping}
\end{equation}
Thus, $q(d_k,t_k)=k$, and each activated layer is uniquely associated
with one charging interval.

The initial layer cannot receive new charging energy:
\begin{equation}
P_{i,0,s,d,t}^{\mathrm{es,ch}}
=
0.
\label{eq:initial_layer_no_charge}
\end{equation}
For each new layer $k\geq1$, charging is permitted only at its associated
formation interval:
\begin{equation}
P_{i,k,s,d,t}^{\mathrm{es,ch}}
=
0,
\qquad
q(d,t)\neq k.
\label{eq:layer_charge_time}
\end{equation}
The activation constraints are
\begin{align}
\varepsilon^{\mathrm{ch}}
u_{i,k,s}^{\mathrm{es,lay}}
&\leq
P_{i,k,s,d_k,t_k}^{\mathrm{es,ch}},
\label{eq:layer_activation_lower}\\
P_{i,k,s,d_k,t_k}^{\mathrm{es,ch}}
&\leq
\overline{P}_i^{\mathrm{es,cap}}
u_{i,k,s}^{\mathrm{es,lay}}.
\label{eq:layer_activation_upper}
\end{align}
A layer cannot discharge before its formation, and the layer variables
are linked to the activation state:
\begin{align}
P_{i,k,s,d,t}^{\mathrm{es,dc}}
&=
0,
\qquad
q(d,t)<k,
\label{eq:no_discharge_before_formation}\\
0
\leq
P_{i,k,s,d,t}^{\mathrm{es,dc}}
&\leq
\overline{P}_i^{\mathrm{es,cap}}
u_{i,k,s}^{\mathrm{es,lay}},
\label{eq:layer_discharge_activation}\\
0
\leq
e_{i,k,s,d,t}^{\mathrm{es}}
&\leq
\overline{E}_i^{\mathrm{es,cap}}
u_{i,k,s}^{\mathrm{es,lay}},
\label{eq:layer_energy_activation}\\
0
\leq
m_{i,k,s,d,t}^{\mathrm{es}}
&\leq
\mathscr{w}^{\max}
\overline{E}_i^{\mathrm{es,cap}}
u_{i,k,s}^{\mathrm{es,lay}}.
\label{eq:layer_carbon_activation}
\end{align}

Variable $u_{i,k,s}^{\mathrm{es,lay}}$ indicates whether the candidate
charging batch is created, and $\varepsilon^{\mathrm{ch}}$ is the minimum
effective charging power. Since each layer receives charging energy only
once, its carbon intensity remains fixed after formation.

\textit{b) Vintage-wise energy and carbon dynamics:}

Define the energy-retention factor as
\begin{equation}
\kappa_i^{\mathrm{es}}
=
1-\sigma_i,
\label{eq:energy_retention}
\end{equation}
where $\sigma_i$ is the self-discharge rate per interval. The layer-wise
energy dynamics are
\begin{align}
e_{i,k,s,d,t+1}^{\mathrm{es}}
=
&
\kappa_i^{\mathrm{es}}
e_{i,k,s,d,t}^{\mathrm{es}}
+
\eta_i^{\mathrm{ch}}
P_{i,k,s,d,t}^{\mathrm{es,ch}}
\Delta t
\nonumber\\
&
-
\frac{
P_{i,k,s,d,t}^{\mathrm{es,dc}}
}{
\eta_i^{\mathrm{dc}}
}
\Delta t,
\qquad
t\in\mathcal{T},
\label{eq:layer_energy_dynamic}
\end{align}
subject to
\begin{equation}
\frac{
P_{i,k,s,d,t}^{\mathrm{es,dc}}
}{
\eta_i^{\mathrm{dc}}
}
\Delta t
\leq
\kappa_i^{\mathrm{es}}
e_{i,k,s,d,t}^{\mathrm{es}}.
\label{eq:available_layer_energy}
\end{equation}
The energy remaining at the end of day $d$ is transferred to day $d+1$
without mixing:
\begin{equation}
e_{i,k,s,d+1,1}^{\mathrm{es}}
=
e_{i,k,s,d,T+1}^{\mathrm{es}},
\quad
d=1,\ldots,D^{\mathrm{blk}}-1.
\label{eq:energy_interday}
\end{equation}

An active layer inherits the nodal carbon intensity at its charging
interval. The conditional equality
$\mathscr{w}_{i,k,s}^{\mathrm{es}}=
\mathscr{w}_{i,s,d_k,t_k}$ is imposed through
\begin{align}
\mathscr{w}_{i,k,s}^{\mathrm{es}}
-
\mathscr{w}_{i,s,d_k,t_k}
&\leq
M^{\mathrm{big}}
\left(
1-u_{i,k,s}^{\mathrm{es,lay}}
\right),
\label{eq:carbon_intensity_bigM_1}\\
\mathscr{w}_{i,s,d_k,t_k}
-
\mathscr{w}_{i,k,s}^{\mathrm{es}}
&\leq
M^{\mathrm{big}}
\left(
1-u_{i,k,s}^{\mathrm{es,lay}}
\right),
\label{eq:carbon_intensity_bigM_2}\\
0
\leq
\mathscr{w}_{i,k,s}^{\mathrm{es}}
&\leq
\mathscr{w}^{\max}
u_{i,k,s}^{\mathrm{es,lay}}.
\label{eq:carbon_intensity_bound}
\end{align}

The carbon inventory of each layer satisfies
\begin{equation}
m_{i,k,s,d,t}^{\mathrm{es}}
=
\mathscr{w}_{i,k,s}^{\mathrm{es}}
e_{i,k,s,d,t}^{\mathrm{es}},
\label{eq:carbon_energy_relation}
\end{equation}
and evolves according to
\begin{align}
m_{i,k,s,d,t+1}^{\mathrm{es}}
=
&
\kappa_i^{\mathrm{es}}
m_{i,k,s,d,t}^{\mathrm{es}}
+
\mathscr{w}_{i,s,d,t}
\eta_i^{\mathrm{ch}}
P_{i,k,s,d,t}^{\mathrm{es,ch}}
\Delta t
\nonumber\\
&
-
\mathscr{w}_{i,k,s}^{\mathrm{es}}
\frac{
P_{i,k,s,d,t}^{\mathrm{es,dc}}
}{
\eta_i^{\mathrm{dc}}
}
\Delta t.
\label{eq:layer_carbon_dynamic}
\end{align}
The layer carbon inventory is transferred between consecutive days as
\begin{equation}
m_{i,k,s,d+1,1}^{\mathrm{es}}
=
m_{i,k,s,d,T+1}^{\mathrm{es}},
\quad
d=1,\ldots,D^{\mathrm{blk}}-1.
\label{eq:carbon_interday}
\end{equation}

In \eqref{eq:layer_energy_dynamic}--\eqref{eq:carbon_interday},
$\eta_i^{\mathrm{ch}}$ and $\eta_i^{\mathrm{dc}}$ are the charging and
discharging efficiencies. The charging term carries the nodal carbon
intensity at the charging interval, whereas the discharging term removes
carbon according to the fixed intensity of the corresponding layer.
Therefore, both energy and carbon attributes are transferred consistently
across the five-day horizon.

\textit{c) Boundary conditions:}

Each scenario is initialized independently. The initial storage inventory
is assigned to layer $k=0$:
\begin{align}
e_{i,0,s,1,1}^{\mathrm{es}}
&=
\mathrm{SOC}_{i,s}^{\mathrm{ini}}
E_i^{\mathrm{es,cap}},
\label{eq:initial_energy}\\
m_{i,0,s,1,1}^{\mathrm{es}}
&=
\mathscr{w}_{i,s}^{\mathrm{es,ini}}
e_{i,0,s,1,1}^{\mathrm{es}},
\label{eq:initial_carbon}\\
\mathscr{w}_{i,0,s}^{\mathrm{es}}
&=
\mathscr{w}_{i,s}^{\mathrm{es,ini}}.
\label{eq:initial_carbon_intensity}
\end{align}
All new layers are initially empty:
\begin{equation}
e_{i,k,s,1,1}^{\mathrm{es}}
=
0,
\qquad
m_{i,k,s,1,1}^{\mathrm{es}}
=
0,
\qquad
k\geq1.
\label{eq:new_layer_initial_state}
\end{equation}

To avoid end-of-horizon depletion, aggregate energy and carbon cycles are
imposed only at the end of the entire five-day block:
\begin{align}
\sum_{k\in\mathcal{K}}
e_{i,k,s,D^{\mathrm{blk}},T+1}^{\mathrm{es}}
&=
e_{i,0,s,1,1}^{\mathrm{es}},
\label{eq:block_energy_cycle}\\
\sum_{k\in\mathcal{K}}
m_{i,k,s,D^{\mathrm{blk}},T+1}^{\mathrm{es}}
&=
m_{i,0,s,1,1}^{\mathrm{es}}.
\label{eq:block_carbon_cycle}
\end{align}

Parameters $\mathrm{SOC}_{i,s}^{\mathrm{ini}}$ and
$\mathscr{w}_{i,s}^{\mathrm{es,ini}}$ denote the initial state of charge
and initial inventory carbon intensity. Different seasonal scenarios are
not linked to each other; hence, the model captures short-term inter-day
transfer rather than seasonal storage transfer. The terminal constraints
restore only the aggregate energy and carbon inventories and do not force
each newly created layer to be emptied.

\emph{4) Dispatchable Generation, PV, and Grid:}

The local-generator output limits are
\begin{equation}
\underline{P}_{g,s,d,t}
\leq
P_{g,s,d,t}
\leq
\overline{P}_{g,s,d,t},
\label{eq:generator_active_limit}
\end{equation}
\begin{equation}
\underline{Q}_{g,s,d,t}
\leq
Q_{g,s,d,t}
\leq
\overline{Q}_{g,s,d,t}.
\label{eq:generator_reactive_limit}
\end{equation}
The initial, intra-day, and inter-day ramping constraints are
\begin{align}
\underline{\Delta}_g
&\leq
P_{g,s,1,1}
-
P_{g,s}^{\mathrm{ini}}
\leq
\overline{\Delta}_g,
\label{eq:initial_ramp}\\
\underline{\Delta}_g
&\leq
P_{g,s,d,t}
-
P_{g,s,d,t-1}
\leq
\overline{\Delta}_g,
\quad
t=2,\ldots,T,
\label{eq:intraday_ramp}\\
\underline{\Delta}_g
&\leq
P_{g,s,d+1,1}
-
P_{g,s,d,T}
\leq
\overline{\Delta}_g,
\quad
d=1,\ldots,D^{\mathrm{blk}}-1.
\label{eq:interday_ramp}
\end{align}

Photovoltaic utilization and curtailment satisfy
\begin{equation}
P_{i,s,d,t}^{\mathrm{PV}}
+
P_{i,s,d,t}^{\mathrm{curt}}
=
\overline{P}_{i,s,d,t}^{\mathrm{PV}},
\label{eq:PV_balance}
\end{equation}
\begin{equation}
P_{i,s,d,t}^{\mathrm{PV}}
\geq
0,
\qquad
P_{i,s,d,t}^{\mathrm{curt}}
\geq
0.
\label{eq:PV_nonnegative}
\end{equation}
The grid-purchase power is limited by
\begin{equation}
0
\leq
P_{s,d,t}^{\mathrm{grid}}
\leq
a_{s,d,t}^{\mathrm{grid}}
\overline{P}^{\mathrm{grid}}.
\label{eq:grid_purchase_limit}
\end{equation}
The annual peak variable satisfies
\begin{equation}
P_{s,d,t}^{\mathrm{grid}}\leq P^{\mathrm{peak}},
\qquad \forall s,d,t.
\label{eq:peak_constraint}
\end{equation}

Parameter $a_{s,d,t}^{\mathrm{grid}}$ is the grid-availability indicator.
Constraint \eqref{eq:interday_ramp} prevents an artificial generator
output jump at a midnight boundary.

\emph{5) LinDistFlow Network Model:}
Each line $(i,j)\in\mathcal{L}$ is oriented from the substation toward
the downstream bus. Let $P_{ij,s,d,t}$ and $Q_{ij,s,d,t}$ denote the
corresponding active and reactive flows; negative $P_{ij,s,d,t}$ denotes
reverse active-power flow. With $\delta_{i0}=1$ at the substation and zero
otherwise, nodal power balance is enforced by
\begin{align}
&
\sum_{g\in\mathcal{G}_i}
P_{g,s,d,t}
+
P_{i,s,d,t}^{\mathrm{PV}}
+
\delta_{i0}
P_{s,d,t}^{\mathrm{grid}}
+
P_{i,s,d,t}^{\mathrm{es,dc}}
\nonumber\\
&
-
P_{i,s,d,t}^{L}
-
P_{i,s,d,t}^{\mathrm{es,ch}}
=
\sum_{j:(i,j)\in\mathcal{L}}P_{ij,s,d,t}
-
\sum_{h:(h,i)\in\mathcal{L}}P_{hi,s,d,t},
\label{eq:nodal_active_balance}\\
&
\sum_{g\in\mathcal{G}_i}
Q_{g,s,d,t}
-
Q_{i,s,d,t}^{L}
=
\sum_{j:(i,j)\in\mathcal{L}}Q_{ij,s,d,t}
-
\sum_{h:(h,i)\in\mathcal{L}}Q_{hi,s,d,t}.
\label{eq:nodal_reactive_balance}
\end{align}
The LinDistFlow voltage equation \cite{Baran1989DistFlow} and operating
limits are
\begin{align}
v_{j,s,d,t}
&=
v_{i,s,d,t}
-2\left(r_{ij}P_{ij,s,d,t}+x_{ij}Q_{ij,s,d,t}\right),
\label{eq:lindistflow_voltage}\\
P_{ij,s,d,t}^{2}+Q_{ij,s,d,t}^{2}
&\leq
\overline{S}_{ij}^{2},
\label{eq:line_limit}\\
\underline{V}_i^2
\leq
v_{i,s,d,t}
&\leq
\overline{V}_i^2,
\label{eq:voltage_limit}\\
v_{0,s,d,t}
&=
V_0^2.
\label{eq:substation_voltage}
\end{align}

For carbon tracing under possible reverse flow, decompose each active
line flow as $P_{ij,s,d,t}=P_{ij,s,d,t}^{+}-P_{ij,s,d,t}^{-}$, where the
two nonnegative components represent the forward and reverse directions.
Their mutual exclusivity is imposed by
\begin{align}
0\leq P_{ij,s,d,t}^{+}
&\leq \overline{S}_{ij}z_{ij,s,d,t}^{\mathrm{f}},
\label{eq:forward_flow}\\
0\leq P_{ij,s,d,t}^{-}
&\leq \overline{S}_{ij}(1-z_{ij,s,d,t}^{\mathrm{f}}),
\label{eq:reverse_flow}\\
P_{ij,s,d,t}
&=P_{ij,s,d,t}^{+}-P_{ij,s,d,t}^{-},
\qquad z_{ij,s,d,t}^{\mathrm{f}}\in\{0,1\}.
\label{eq:flow_decomposition}
\end{align}

\emph{6) Nodal Carbon Flow and Carbon-Intensity Limits:}
The carbon flow injected by storage discharge is
\begin{equation}
R_{i,s,d,t}^{\mathrm{es,dc}}
=
\sum_{k\in\mathcal{K}}
\mathscr{w}_{i,k,s}^{\mathrm{es}}
P_{i,k,s,d,t}^{\mathrm{es,dc}}.
\label{eq:storage_carbon_flow}
\end{equation}
Under the proportional-sharing principle \cite{chen2025copf}, the nodal
carbon-flow balance is
\begin{align}
&
\mathscr{w}_{i,s,d,t}
\left(
\sum_{g\in\mathcal{G}_i}
P_{g,s,d,t}
+
P_{i,s,d,t}^{\mathrm{PV}}
+
\delta_{i0}
P_{s,d,t}^{\mathrm{grid}}
\right.
\nonumber\\
&
\left.
+
P_{i,s,d,t}^{\mathrm{es,dc}}
+
\sum_{h:(h,i)\in\mathcal{L}}P_{hi,s,d,t}^{+}
+
\sum_{j:(i,j)\in\mathcal{L}}P_{ij,s,d,t}^{-}
\right)
\nonumber\\
=\;&
\sum_{g\in\mathcal{G}_i}
\mathscr{w}_{i,g,s,d,t}^{G}
P_{g,s,d,t}
+
\delta_{i0}
\mathscr{w}_{s,d,t}^{\mathrm{grid}}
P_{s,d,t}^{\mathrm{grid}}
\nonumber\\
&
+
R_{i,s,d,t}^{\mathrm{es,dc}}
+
\sum_{h:(h,i)\in\mathcal{L}}
\mathscr{w}_{h,s,d,t}P_{hi,s,d,t}^{+}
+
\sum_{j:(i,j)\in\mathcal{L}}
\mathscr{w}_{j,s,d,t}P_{ij,s,d,t}^{-}.
\label{eq:nodal_carbon_balance}
\end{align}
The nodal carbon-intensity caps are
\begin{equation}
\mathscr{w}_{i,s,d,t}
\leq
\overline{\mathscr{w}}_i,
\qquad
\forall i,s,d,t.
\label{eq:nodal_carbon_cap}
\end{equation}

In \eqref{eq:nodal_carbon_balance}, the left-hand side is the nodal
carbon intensity multiplied by the total incoming power. The right-hand
side contains carbon flows from local generation, the upstream grid,
storage discharge, and neighboring buses. Photovoltaic generation is
assumed to have zero operational carbon intensity; therefore, its power
appears in the total incoming power but contributes no carbon flow.
Parameter $\overline{\mathscr{w}}_i$ is a bus-specific hard cap and no
carbon-intensity slack variable is introduced.

For interpretation, define the incoming power and carbon flow excluding
storage discharge as
\begin{align}
A_{i,s,d,t}^{C}
=
&
\sum_{g\in\mathcal{G}_i}
P_{g,s,d,t}
+
P_{i,s,d,t}^{\mathrm{PV}}
+
\delta_{i0}
P_{s,d,t}^{\mathrm{grid}}
\nonumber\\
&
+
\sum_{h:(h,i)\in\mathcal{L}}P_{hi,s,d,t}^{+}
+
\sum_{j:(i,j)\in\mathcal{L}}P_{ij,s,d,t}^{-},
\label{eq:other_input_power}\\
B_{i,s,d,t}^{C}
=
&
\sum_{g\in\mathcal{G}_i}
\mathscr{w}_{i,g,s,d,t}^{G}
P_{g,s,d,t}
+
\delta_{i0}
\mathscr{w}_{s,d,t}^{\mathrm{grid}}
P_{s,d,t}^{\mathrm{grid}}
\nonumber\\
&
+
\sum_{h:(h,i)\in\mathcal{L}}
\mathscr{w}_{h,s,d,t}P_{hi,s,d,t}^{+}
+
\sum_{j:(i,j)\in\mathcal{L}}
\mathscr{w}_{j,s,d,t}P_{ij,s,d,t}^{-}.
\label{eq:other_input_carbon}
\end{align}
Combining \eqref{eq:nodal_carbon_balance} and
\eqref{eq:nodal_carbon_cap} yields
\begin{equation}
\sum_{k\in\mathcal{K}}
\left(
\overline{\mathscr{w}}_i
-
\mathscr{w}_{i,k,s}^{\mathrm{es}}
\right)
P_{i,k,s,d,t}^{\mathrm{es,dc}}
\geq
D_{i,s,d,t}^{C},
\label{eq:carbon_gap_constraint}
\end{equation}
where
\begin{equation}
D_{i,s,d,t}^{C}
=
B_{i,s,d,t}^{C}
-
\overline{\mathscr{w}}_i
A_{i,s,d,t}^{C}.
\label{eq:carbon_gap}
\end{equation}

A positive $D_{i,s,d,t}^{C}$ indicates that storage or other low-carbon
resources are required to satisfy the nodal cap. The carbon-intensity
reduction capability per unit discharge of layer $k$ is
$\overline{\mathscr{w}}_i-\mathscr{w}_{i,k,s}^{\mathrm{es}}$. Hence, a
lower-carbon layer provides a larger reduction capability. If the
available low-carbon layers are insufficient, the first-stage model must
increase storage power or energy capacity, alter the installation
location, increase charging during preceding low-carbon intervals, or
adjust the schedules of other resources.

\emph{7) Complete Model:}
Using the decision vectors introduced above, the primary problem is
\begin{equation}
\begin{aligned}
\min_{\boldsymbol{x}^{\mathrm{pl}},\boldsymbol{y}^{\mathrm{op}}}
\quad&
C^{\mathrm{ann}}
\\
\mathrm{s.t.}
\quad&
(\boldsymbol{x}^{\mathrm{pl}},\boldsymbol{y}^{\mathrm{op}})
\in
\Omega,
\end{aligned}
\label{eq:compact_stage1}
\end{equation}
where $\Omega$ is the feasible set defined by
\eqref{eq:peak_constraint}, \eqref{eq:backlog_dynamic}--\eqref{eq:DC_shedding_limit},
\eqref{eq:storage_power_capacity}--\eqref{eq:block_carbon_cycle},
\eqref{eq:generator_active_limit}--\eqref{eq:grid_purchase_limit}, and
\eqref{eq:nodal_active_balance}--\eqref{eq:nodal_carbon_cap}.

The second-stage model is
\begin{equation}
\begin{aligned}
\min_{\boldsymbol{y}^{\mathrm{op}}}
\quad&
J^{\mathrm{res}}
\\
\mathrm{s.t.}
\quad&
\eqref{eq:fix_siting}
-
\eqref{eq:economic_equivalence},
\\
&
(\boldsymbol{x}^{\mathrm{pl},*},\boldsymbol{y}^{\mathrm{op}})
\in
\Omega.
\end{aligned}
\label{eq:compact_stage2}
\end{equation}

Both stages are mixed-integer nonlinear programs because of the operating
binaries and bilinear carbon-accounting terms. Fixing the operating
binaries after Stage I reduces Stage II to a nonlinear program.


\begin{thebibliography}{2}

\bibitem{Baran1989DistFlow}
M.~E. Baran and F.~F. Wu, ``Optimal sizing of capacitors placed on a
radial distribution system,'' \emph{IEEE Trans. Power Del.}, vol.~4,
no.~1, pp.~735--743, Jan.~1989.

\bibitem{chen2025copf}
X.~Chen, A.~Sun, W.~Shi, and N.~Li, ``Carbon-aware optimal power flow,'' \emph{IEEE Trans. Power Syst.}, vol.~40, no.~4, pp.~3090--3104, Jul.~2025.

\end{thebibliography}

\end{document}
